\documentclass[12pt,a4paper]{article}
\usepackage[utf8]{inputenc}
\usepackage[spanish]{babel}
\usepackage{geometry}
\usepackage{enumitem}
\usepackage{titlesec}
\usepackage{xcolor}
\usepackage[hidelinks]{hyperref}
\usepackage{listings}

\geometry{a4paper, margin=1in}

\titleformat{\section}{\normalfont\Large\bfseries\color{blue}}{}{0pt}{}
\titleformat{\subsection}{\normalfont\large\bfseries\color{darkgray}}{}{0pt}{}

\title{Base Teórica - Semana 1 \\ \large Curso: Seguridad Informática}
\author{Prof. César Rodríguez}
\date{\today}

\begin{document}

\maketitle

\tableofcontents
\newpage

\section*{Introducción}
\addcontentsline{toc}{section}{Introducción}
Este documento proporciona los fundamentos teóricos necesarios para desarrollar los módulos correspondientes a la \textbf{Semana 1}. Está estructurado según las responsabilidades de cada grupo de trabajo.

\section{Grupo 1: Reconocimiento DNS}
El Sistema de Nombres de Dominio (DNS) es la agenda telefónica de Internet. Traduce nombres de dominio legibles por humanos (como \texttt{google.com}) a direcciones IP legibles por máquinas.

\subsection{Tipos de Registros DNS Clave}
\begin{itemize}[label=$\bullet$]
    \item \textbf{Registro A (Address):} Asocia un nombre de dominio con una dirección IPv4 (ej. 192.168.1.1). Es el registro más fundamental.
    \item \textbf{Registro AAAA:} Funciona exactamente igual que el registro A, pero asocia el dominio a una dirección IPv6.
    \item \textbf{Registro MX (Mail Exchange):} Indica qué servidores de correo están autorizados a aceptar correos electrónicos en nombre del dominio.
    \item \textbf{Registro NS (Name Server):} Delega una zona DNS para utilizar los servidores de nombres especificados. Indica dónde se administran realmente los registros del dominio.
    \item \textbf{Registro TXT (Text):} Permite a los administradores insertar texto arbitrario en el registro DNS. Se usa comúnmente para verificación de dominio y políticas de seguridad de correo (como SPF, DKIM o DMARC).
    \item \textbf{Registro SOA (Start of Authority):} Contiene información administrativa sobre la zona DNS, como el correo del administrador, el número de serie de la zona y los tiempos de actualización.
\end{itemize}

\textit{Sugerencia de implementación:} La librería \texttt{dnspython} proporciona el método \texttt{dns.resolver.resolve()} que facilita la consulta de cualquiera de estos registros.

\section{Grupo 2: OSINT y Huella Digital}
OSINT (Open-Source Intelligence) es la recopilación y análisis de información obtenida de fuentes públicas y abiertas.

\subsection{WHOIS}
WHOIS es un protocolo de consulta y respuesta ampliamente utilizado para consultar bases de datos que almacenan los usuarios registrados o asignatarios de un recurso de Internet, como un nombre de dominio o un bloque de direcciones IP. 
\begin{itemize}
    \item Proporciona: Fechas de creación/expiración, registrador, correos electrónicos de contacto y servidores DNS asociados.
\end{itemize}
\textit{Sugerencia:} En Python, la librería \texttt{python-whois} facilita esta extracción devolviendo un diccionario manejable.

\subsection{Google Dorks (Google Hacking)}
Consiste en utilizar operadores de búsqueda avanzados de Google para encontrar información sensible que ha sido indexada por el buscador (a menudo por error de configuración del administrador).
\begin{itemize}
    \item \texttt{site:dominio.com} - Restringe los resultados a un dominio específico.
    \item \texttt{-www} - Excluye el subdominio 'www', revelando a menudo otros subdominios menos conocidos (ej. \texttt{site:ejemplo.com -www}).
    \item \texttt{filetype:ext} - Busca archivos con una extensión específica (ej. \texttt{filetype:log}, \texttt{filetype:env}, \texttt{filetype:sql}).
    \item \texttt{intitle:} o \texttt{inurl:} - Busca cadenas específicas en el título de la página o en la URL (ej. \texttt{intitle:"index of"}).
\end{itemize}

\section{Grupo 3: Descubrimiento de Hosts}
Antes de atacar o auditar, se debe saber qué equipos están "vivos" (encendidos y conectados a la red).

\subsection{Ping Sweep (ICMP Echo Request)}
Es la técnica más básica. Envía un paquete \texttt{ICMP Type 8} (Echo Request) a cada IP en una red. Si el host está vivo y no bloquea ICMP, responde con un \texttt{ICMP Type 0} (Echo Reply).

\subsection{Rangos de Red y CIDR}
CIDR (Classless Inter-Domain Routing) es una forma de representar rangos de IPs. Por ejemplo, \texttt{192.168.1.0/24} representa las 256 direcciones desde \texttt{192.168.1.0} hasta \texttt{192.168.1.255}.
\textit{Sugerencia:} Python incluye el módulo \texttt{ipaddress} que facilita iterar sobre rangos de red usando notación CIDR (\texttt{ipaddress.ip\_network()}).

\subsection{TCP SYN Ping}
Muchos firewalls modernos bloquean el tráfico ICMP (ping tradicional). Para eludirlos, los atacantes y auditores envían un paquete TCP con el flag SYN (intento de conexión) a un puerto que suele estar abierto (como el 80 o el 443).
\begin{itemize}
    \item Si el puerto está abierto, el host responde con SYN-ACK.
    \item Si el puerto está cerrado, responde con RST (Reset).
    \item En ambos casos, el auditor confirma que el host está \textbf{vivo}.
\end{itemize}

\section{Grupo 4: Escaneo de Puertos y Servicios}
Una vez descubierto un host, el siguiente paso es descubrir qué servicios (y en qué puertos) está ejecutando.

\subsection{Escáner TCP Connect}
Se basa en completar el "Three-way Handshake" (saludo de tres vías) completo de TCP:
\begin{enumerate}
    \item El auditor envía un paquete \textbf{SYN}.
    \item El objetivo responde con \textbf{SYN-ACK} (si el puerto está abierto).
    \item El auditor responde con \textbf{ACK}, completando la conexión, y luego la cierra.
\end{enumerate}
Es el escaneo más sencillo de programar en Python usando la librería estándar \texttt{socket} (\texttt{socket.connect\_ex()}), aunque es ruidoso y queda registrado en los logs del servidor.

\subsection{Concurrencia}
Escanear 65,535 puertos uno por uno de forma secuencial tomaría horas. Para acelerar el proceso, se usan \textbf{hilos (threads)}. Un hilo permite lanzar múltiples intentos de conexión de manera casi simultánea. En Python se puede usar \texttt{threading} o, mejor aún, \texttt{concurrent.futures.ThreadPoolExecutor}.

\subsection{Escaneo UDP}
A diferencia de TCP, UDP no tiene estado y no realiza un saludo inicial.
\begin{itemize}
    \item Se envía un paquete UDP vacío al puerto.
    \item Si no hay respuesta, el puerto suele estar \textbf{abierto} o filtrado (el servicio recibió el paquete y lo ignoró al no entenderlo).
    \item Si el puerto está cerrado, el host normalmente responde con un paquete \textbf{ICMP Port Unreachable} (Tipo 3, Código 3).
\end{itemize}
Dado que los paquetes se pueden perder en la red, el escaneo UDP suele ser más lento e impreciso que el TCP.

\subsection{Banner Grabbing}
Saber que el puerto 22 o el 80 están abiertos no es suficiente. El "Banner Grabbing" consiste en conectarse al puerto abierto, enviar unos bytes de prueba (o simplemente esperar a que el servidor hable) y capturar el texto que el servicio devuelve de forma predeterminada (el "Banner"). Este banner suele revelar el nombre y la versión del software (ej. \texttt{OpenSSH 8.2p1} o \texttt{Apache/2.4.41}), lo cual es fundamental para buscar vulnerabilidades asociadas a esa versión.

\newpage
\section*{Glosario de Términos}
\addcontentsline{toc}{section}{Glosario de Términos}
\begin{description}
    \item[Banner Grabbing] Técnica de recolección de información que consiste en conectarse a un puerto abierto y leer el mensaje de bienvenida (``banner'') para identificar el servicio y su versión.
    \item[CIDR] \textit{Classless Inter-Domain Routing} (Enrutamiento Inter-Dominios sin Clases). Método para asignar direcciones IP y enrutar paquetes, comúnmente usado para denotar rangos de red (ej. /24).
    \item[DNS] \textit{Domain Name System} (Sistema de Nombres de Dominio). Servicio que traduce nombres legibles (dominios) en direcciones IP numéricas.
    \item[Google Dorks] Uso de operadores de búsqueda avanzada en Google para encontrar información sensible, archivos expuestos o vulnerabilidades no intencionadas.
    \item[ICMP] \textit{Internet Control Message Protocol}. Protocolo de red utilizado para enviar mensajes de error e información operativa, como el comando de red \texttt{ping}.
    \item[OSINT] \textit{Open-Source Intelligence} (Inteligencia de Fuentes Abiertas). Recopilación y análisis de información de fuentes públicas disponibles.
    \item[TCP] \textit{Transmission Control Protocol}. Protocolo de comunicaciones orientado a la conexión que garantiza la entrega ordenada y sin errores de los paquetes de datos.
    \item[UDP] \textit{User Datagram Protocol}. Protocolo de comunicaciones sin conexión que envía paquetes sin garantizar su entrega ni orden, siendo más rápido pero menos fiable que TCP.
    \item[WHOIS] Protocolo de red utilizado para consultar bases de datos públicas sobre el registro de nombres de dominio y bloques de direcciones IP.
\end{description}

\end{document}